\begin{figure}
    \small
	\centering
	\begin{align*}
		\decompose(e : \tau) =
		\begin{cases}
			(x : \tau, ())
				& \text{if } e = x : \tau \\[8pt]
			(\true : \bool, ())
				& \text{if } e = \true : \bool \\[8pt]
			(\false : \bool, ())
				& \text{if } e = \false : \bool \\[8pt]
			(\finconst{k}{n} : \fin{n}, ())
				& \text{if } e = \finconst{k}{n} : \fin{n} \\[8pt]
			(\diverge : \tau, ())
				& \text{if } e = \diverge : \tau \\[8pt]
			(() : \unit, ())
				& \text{if } e = () : \unit \\[8pt]
			((\square, \floattype{B}{V}) : \floattype{B}{V}, c)
				& \text{if } e = c : \floattype{B}{V} \\[8pt]
			(\discrete(p_0,\ldots,p_{n-1}) : \fin{n}, ())
				& \text{if } e = \discrete(p_0,\ldots,p_{n-1}) : \fin{n} \\[8pt]
			(\uniform(\decompose(e_1)_1, \decompose(e_2)_1) : \floattype{B}{V}, \\
			(\decompose(e_1)_2, \decompose(e_2)_2))
				& \text{if } e = \uniform(e_1, e_2) : \floattype{B}{V} \\[8pt]
			(\decompose(e_1)_1 \sim \decompose(e_2)_1 : \bool, \\
			(\decompose(e_1)_2, \decompose(e_2)_2))
				& \text{if } e = e_1 \sim e_2 : \bool, \sim = <, \finlt{n}, \fineq{n} \\[8pt]
            (\ifkw\; \decompose(e_1)_1\; \thenkw\; \decompose(e_2)_1\; \elsekw\; \decompose(e_3)_1 : \tau, \\
            (\decompose(e_1)_2, \decompose(e_2)_2, \decompose(e_3)_2))
                & \text{if } e = \ifkw\; e_1\; \thenkw\; e_2\; \elsekw\; e_3 : \tau \\[8pt]
            (\letkw\; x = \decompose(e_1)_1\; \inkw\; \decompose(e_2)_1 : \tau, \\
            (\decompose(e_1)_2, \decompose(e_2)_2))
                & \text{if } e = \letkw\; x = e_1\; \inkw\; e_2 : \tau \\[8pt]
            (\observekw(\decompose(e_1)_1 : \bool) : \unit, \decompose(e_1)_2)
                & \text{if } e = \observekw(e_1 : \bool) : \unit \\[8pt]
		\end{cases}
	\end{align*}
	\caption{Recursive definition of $\decompose$, a decomposition function, over typed expressions in $Expr^{\mathsf U+}$. $\decompose(\cdot)_1$ and $\decompose(\cdot)_2$ denote first and second projections of $\decompose$, respectively. $\decompose(\cdot)_2$ is canonically bijective to a flat vector in $\R^n$ via coordinate concatenation over holes.}
	\label{fig:decompose-definition-1}
\end{figure}

\begin{figure}
    \small
	\centering
	\begin{align*}
		\decompose(e : \tau) =
		\begin{cases}
            (\decompose(e_1)_1 \logand \decompose(e_2)_1 : \bool, \\
			(\decompose(e_1)_2, \decompose(e_2)_2))
				& \text{if } e = e_1 \logand e_2 : \bool \\[8pt]
			(\decompose(e_1)_1 \logor \decompose(e_2)_1 : \bool, \\
			(\decompose(e_1)_2, \decompose(e_2)_2))
				& \text{if } e = e_1 \logor e_2 : \bool \\[8pt]
			(\lognot\; \decompose(e_1)_1 : \bool, \decompose(e_1)_2)
				& \text{if } e = \lognot\; e_1 : \bool \\[8pt]
			((\decompose(e_1)_1, \decompose(e_2)_1) : \tau_1 * \tau_2, \\
			(\decompose(e_1)_2, \decompose(e_2)_2))
				& \text{if } e = (e_1, e_2) : \tau_1 * \tau_2 \\[8pt]
			(\fstkw\; \decompose(e_1)_1 : \tau_1, \decompose(e_1)_2)
				& \text{if } e = \fstkw\; e_1 : \tau_1 \\[8pt]
			(\sndkw\; \decompose(e_1)_1 : \tau_2, \decompose(e_1)_2)
				& \text{if } e = \sndkw\; e_1 : \tau_2 \\[8pt]
			(\funkw\; x \to \decompose(e_1)_1 : \tau_1 \to \tau_2, \decompose(e_1)_2)
				& \text{if } e = \funkw\; x \to e_1 : \tau_1 \to \tau_2 \\[8pt]
			(\fixkw\; f\; x := \decompose(e_1)_1 : \tau_1 \to \tau_2, \decompose(e_1)_2)
				& \text{if } e = \fixkw\; f\; x := e_1 : \tau_1 \to \tau_2 \\[8pt]
			(\decompose(e_1)_1\; \decompose(e_2)_1 : \tau_2, \\
			(\decompose(e_1)_2, \decompose(e_2)_2))
				& \text{if } e = (e_1 : \tau_1 \to \tau_2\;\; e_2 : \tau_1) : \tau_2 \\[8pt]
			(\textnormal{\ttfamily nil} : \listty{\tau_1}, ())
				& \text{if } e = \textnormal{\ttfamily nil} : \listty{\tau_1} \\[8pt]
			(\decompose(e_1)_1 :: \decompose(e_2)_1 : \listty{\tau_1}, \\
			(\decompose(e_1)_2, \decompose(e_2)_2))
				& \text{if } e = e_1 :: e_2 : \listty{\tau_1} \\[8pt]
			(\matchkw\; \decompose(e_0)_1\; \withkw\; \matchcase \textnormal{\ttfamily nil} \rightarrow \decompose(e_{\textnormal{nil}})_1 \\
			\matchcase h::t \rightarrow \decompose(e_{\textnormal{cons}})_1\; {\textnormal{\ttfamily\bfseries end}} : \tau_2, 
                & \text{if } e = (\matchkw\; e_0 : \listty{\tau_1}\; \withkw\; \matchcase \textnormal{\ttfamily nil} \rightarrow e_{\textnormal{nil}} : \tau_2 \\
			(\decompose(e_0)_2, \decompose(e_{\textnormal{nil}})_2, \decompose(e_{\textnormal{cons}})_2))
                &\qquad\qquad\;\;\;\matchcase h::t \rightarrow e_{\textnormal{cons}} : \tau_2\; {\textnormal{\ttfamily\bfseries end}}) : \tau_2
		\end{cases}
	\end{align*}
	\caption{Recursive definition of $\decompose$ over typed expressions in $Expr^{\mathsf U+}$ (continued).}
	\label{fig:decompose-definition-2}
\end{figure}
